\documentclass{article}
\usepackage{graphicx}
\usepackage{float}
\usepackage{amsmath}
\usepackage{subcaption}
\usepackage[dvipsnames]{xcolor}

\title{Globalization and Structural Change in the Phillips Curve}
\author{Brendan Stimola}
\date{May 2026}

\begin{document}

\maketitle

\section{Abstract}

This paper studies whether globalization and China joining the WTO in particular weakens the Phillips Curve relationship in the United States, and how the Phillips Curve evolved over time using different measures for inflation. The first model includes unemployment, trade openness, interaction effects, and lagged inflation. Ordinary Least Squares (OLS) regression is used to estimate the inflation dynamics. The second model uses China's accession into the WTO for a structural break. The model uses 4 variables: unemployment, Post WTO (the immediate drop in inflation following China joining the WTO) and POST UNRATE (The change in slope of the Phillips Curve after China joined the WTO) and lagged inflation. To measure a structural break with the second model there was a Wald test to check for a structural break after 2001.The final model replaces inflation with the inflation gap ($\pi-\overline{\pi}$).The results suggest that inflation is highly persistent and that globalization weakens the relationship between unemployment and inflation. The Phillips Curve remains present in the United States, but globalization weakens the relationship between unemployment and inflation. However, once the inflation gap is accounted for, the globalization effect becomes substantially weaker, suggesting that globalization primarily influences inflation through inflation persistence. Finally in one of the models all of the primary variables of importance, had p-values above 0.5 rendering them useless. These findings are consistent with theories that the Phillips Curve has structurally changed after China joined the WTO due to a more globalized economy. 

\section{Introduction}

The Phillips Curve is the theory that unemployment and inflation are negatively correlated. Modern-day economists care because this relationship has ripple effects everywhere, from Fed policy to inflation control and employment. The present issue is that the relationship may have weakened due to a variety of reasons. Most of them have to do with how globalization affects prices in an economy. The research question is does increasing trade openness, specifically after China joined the WTO. Change how powerful the Phillips Curve is? The baseline model utilizes an Ordinary Least Squares (OLS) framework to isolate the relationship between trade to GDP and the unemployment rate. The second model uses China joining the WTO, which is a pivotal moment in globalization to see the effect on inflation and unemployment. The third model does the same thing as the first excepts it replaces inflation with the gap between inflation and average inflation.

\section{Literature Review}

Modern day studies say that globalization is flattening the Phillips Curve (International Journal of Central Banking). The IJCB as its also known attributes this to decoupling demand and supply from domestic markets. Some also state how the output gaps effect on inflation is shrinking because of globalization. Causing pay and unemployment to be based less on the domestic economy (Science Direct). The "Becker Friedman Institute" studies how regional Phillips Curves hold up, and came to the conclusion that 40 percent of the decline in correlation is due to imports (Becker Friedman Institute).

\section{Economic Background}

Inflation is the year over year change in the CPI. The CPI uses a basket of goods and services to see the level of the price in an economy. Unemployment is the percent of the labor force that is looking for work but can not find it. Trade openness is measured as:
$$(\frac{Imports+Exports}{GDP})*100$$
This measures the amount of economic activity that is due to trade. This is used as a proxy for globalization as it measures the percent of the economy that is reliant on trade. Lagged inflation is inflation shifted one period. This is used to see how much of today's inflation is depended on the inflation of the past. This is important because it shows how inflation expectations affect inflation in the long-term. The equation for the Phillips Curve is 
$$
\pi_t=\pi^e_t-a(u_t-u_n)+v_t
$$
$\pi^e_t$ are the current expectations of inflation. $\alpha$ is the sensitivity of the Phillips Curve, $u_t$ and $u_n$ are the current unemployment and the natural rate of unemployment respectively; And $v_t$ are shocks to the economy (oil shocks etc..). China joining the WTO was a very important moment in history as it marked the beginning of an influx of low-cost imports from China due to increased supply and lower domestic demand to the USA.
\section{Data}

For inflation in model 1 the year over year change in the CPI was used. For inflation in model 3 the inflation (from model 1) - average inflation was calculated to get the inflation gap. For the unemployment source the domestic unemployment rate is measured to get a picture of the labor market. For the trade term trade to GDP was picked to capture the openness of the US economy. Lagged Inflation was introduced to capture the momentum and persistence of inflation. In the model there is a "dummy variable" to split the regression into pre- and post-2001. All of these economic terms are known as CPIAUCSL(inflation), UNRATE(unemployment rate), EXPGS and IMPGS (added together are trade), GDP(size of the economy). This was chosen as they are standard macro variables used in assessing the Phillips Curve. Below is a figure of trade to GDP which is imports plus exports over GDP.

\begin{figure}[H] 
\centering 
\includegraphics[width=0.85\textwidth]{trade_GDP.png} 
\caption{A figure showing the trade to GDP, a common measure of globalization, from 1954 to 2026 in the United States.} 
\label{fig:trade_GDP_line} 
\end{figure}

\section{Methodology}

For the regression method OLS was chosen. The purpose was the analyze the relationship between inflation, the unemployment rate and how globalized an economy is. The interaction term tests to what degree globalization effects the Phillips Curve. The equation for the first model is:
\begin{align*}
\text{Inflation}_t =
& \ \beta_0 + \beta_1 \text{Unemployment}_t + \beta_2 \text{Trade}_t \\
& + \beta_3 \text{Interaction}_t + \beta_4 \text{Inflation}_{t-1} + \varepsilon_t
\end{align*}
The equation for the second model is:
\begin{align*}
\text{Inflation}_t =
& \ \beta_0 + \beta_1 \text{Unemployment}_t + \beta_2 \text{Post UNRATE}_t \\
& + \beta_3 \text{Post WTO}_t + \text{Inflation}_{t-1} + \varepsilon_t
\end{align*}
The equation for the third model is:
\begin{align*}
\text{Inflation Gap}_t =
& \ \beta_0 + \beta_1 \text{Unemployment}_t + \beta_2 \text{Trade}_t \\
& + \beta_3 \text{Interaction}_t + \text{Inflation}_{t-1} + \varepsilon_t
\end{align*}
Where $\pi_t$ is the current inflation. This equation uses the betas as coefficients to measure the line of best fit. Each variable is multiplied by its beta and then the error is added to try to estimate the inflation or inflation gap ($\pi_t-\overline{\pi}$). The average inflation in the inflation gap, is the average of inflation in the US from 1954-2026.
\section{Results}

\begin{table}[H]
\centering
\begin{tabular}{|c|c|c|c|}
\hline
Variable & Coefficient & P-value & Std Err \\
\hline
UNRATE & -0.267 & $<0.01$ & 0.067 \\
TRADE & 0.012 & 0.296 & 0.011\\
Interaction & 0.023 & $<0.01$ & 0.008 \\
Lagged Inflation & 1.017 & $<0.01$ & 0.016 \\
\hline
\end{tabular}
\caption{Regression Results for Model 1}
\end{table}

\begin{table}[H]
\centering
\begin{tabular}{|c|c|c|c|}
\hline
Variable & Coefficient & P-value & Std Err \\
\hline
UNRATE & -0.059 & 0.558 & 0.100\\
Post WTO & 0.065 & 0.730 & 0.188\\
Post UNRATE & -0.060 & 0.599 & 0.115\\
Lagged Inflation & 0.946 & $<0.01$ & 0.035\\
\hline
\end{tabular}
\caption{Regression Results for Model 2 (lagged inflation)}
\end{table}

\begin{table}[H]
\centering
\begin{tabular}{|c|c|c|c|}
\hline
Variable & Coefficient & P-value & Std Err \\
\hline
UNRATE & -0.267 &    $<0.01$ & 0.330\\
Post WTO &  -2.191  &  0.096  &  1.315\\
Post UNRATE &  -2.731  &   $<0.01$  & 0.534\\
\hline
\end{tabular}
\caption{Regression Results for Model 2 (without lagged inflation)}
\end{table}

\begin{table}[H]
\centering
\begin{tabular}{|c|c|c|c|}
\hline
Variable & Coefficient & P-value & Std Err \\
\hline
UNRATE & -0.107 & 0.01 & 0.041\\
TRADE & 0.037 & 0.239 & 0.031\\
Interaction & -0.002 & 0.848 & 0.013\\
Lagged Inflation & 0.961 & $<0.01$ & 0.030\\
\hline
\end{tabular}
\caption{Regression Results for Model 3}
\end{table}

\begin{table}[H]
\centering
\begin{tabular}{|c|c|c|c|}
\hline
Model & $R^2$ & Durbin-Watson & Cond. No. \\
\hline
Model 1 & 0.983 & 1.187 & 28.2\\
Model 2 & 0.940 & 1.478 & 40.5\\
Model 3 & 0.940 & 1.503 & 6.74\\
\hline
\end{tabular}
\caption{Results for the models}
\end{table}

The results for the final table suggest that most of the data is overall correct but some p-values are too high to be used. Adding lagged inflation increases the r squared and the Durbin-Watson, but also removes the noise that the variables used to absorb leaving a more cyclical pattern that raises p-values. Overall though besides the second model which used a different model and caused higher unusable p-values. The models overall have pretty good results and can be trusted.

\subsection{The normal Phillips Curve (with trade to GDP and interaction term) from 1954 to 2026}
The unemployment is negative and statistically significant, meaningly that the Phillips Curve does still exist but maybe just weaker. The lagged inflation is very close to 1 and has a p-value less than 0.05 meaning inflation is very dependent on inflation from the past. The trade term is positive but not statistically significant (P = 0.296), so we fail to reject the null hypothesis that trade openness has an effect on inflation. The trade times unemployment rate though yielded economically meaningful results. The interaction term is positive meaning globalization does flatten the Phillips Curve, and the p-value is less than 0.01 meaning it is statistically significant. 
\subsection{The before and after China joined WTO}
In model 2 there are some interesting results, even though they may be hard to see. In model 2 with lagged inflation only Lagged Inflation is significant. It has a p value of less than 0.01 and a coefficient of 0.961. While not visible in just the table, this is important data to see. When there is no lagged inflation the variables absorb the error terms. Causing their p-values to inflate and make the model look significant. When adding the lagged inflation, all that remains is the cyclical cycle underneath. These variables not being significant proves that China joining the WTO may have caused a structural break, but the cyclical variables never changed. This pattern is consistent with what can be observed when the lagged inflation is removed. Looking at the full picture not just cyclical, it is clear that there may have been a structural break. The Wald test came back positive signifying that China joining the WTO structurally changed the US economy.
\subsection{Inflation gap model}
However, In model 3 some of the variables are statistically insignificant. The reason for this is once the inflation gap is accounted for all that is left is the short-term trends that makes unemployment harder to track. Once the long-term movement is removed the pattern is very cyclical and could cause the model to not function and give results with weaker explanatory power. The unemployment rate is still significant with a negative coefficient and significant p-value. Just like the lagged inflation which also acted just like it did in the other model.

\subsection{Structural Change in the Phillips Curve}

\begin{figure}[H]
    \centering
    \includegraphics[width=0.85\textwidth]
    {1954_2001.png}

    \caption{
    A plot of the relationship between unemployment and inflation between 1954 and 2001. The 3 lines are the Phillips curve during different percentiles of globalization in that period. If the line is positive it is the opposite of the Phillips Curve and vice versa. The steeper the line the more sensitive inflation is to unemployment.
    }

    \label{fig:1954_2001}
\end{figure}

\begin{figure}[H]
    \centering
    \includegraphics[width=0.85\textwidth]
    {2001_2026.png}

    \caption{
    A plot of the relationship between unemployment and inflation between 2001 and 2026. The 3 lines are the Phillips curve during different percentiles of globalization in that period. If the line is positive it is the opposite of the Phillips Curve and vice versa. The steeper the line the more sensitive inflation is to unemployment.
    }

    \label{fig:2001_2026}
\end{figure}

Above are 2 figures are showing the Phillips Curve relationship in two periods. 1954-2001 (before China joined the WTO) and 2001-2026 (after China joined the WTO). A crucial nuance is observable in figure two,  it may seem that high globalization is just like low globalization and globalization has opposite effect on Phillips Curve. Because globalization is high after 2001 low globalization is still high by historical standards. Also because of the recessions where inflation drops and unemployment rises the data could be slightly unreliable.

\section{Diagnostics and Limitations}
\subsection{Diagnostics} 
To verify the core OLS assumptions of normality and linearity, I evaluated three diagnostic plots: a Quantile-Quantile (Q-Q) plot, a residuals histogram, and a residuals vs. fitted values plot (Figure \ref{fig:diagnostics}). 

\begin{figure}[H]
     \centering
     \begin{subfigure}[b]{0.32\textwidth}
         \centering
         \includegraphics[width=\textwidth]{QQ_PLOT.png}
         \caption{Q-Q Plot}
         \label{fig:qq_plot}
     \end{subfigure}
     \hfill
     \begin{subfigure}[b]{0.32\textwidth}
         \centering
         \includegraphics[width=\textwidth]{Residuals_histogram.png}
         \caption{Residuals Histogram}
         \label{fig:residuals_histogram}
     \end{subfigure}
     \hfill
     \begin{subfigure}[b]{0.32\textwidth}
         \centering
         \includegraphics[width=\textwidth]{Residuals_vs_Fitted.png}
         \caption{Residuals vs. Fitted}
         \label{fig:residuals_vs_fitted}
     \end{subfigure}
        \caption{Model diagnostic plots. The Q-Q plot (a) shows residuals aligning with the 45-degree normal reference line. The histogram (b) confirms a symmetric, bell-shaped error distribution. The residuals vs. fitted plot (c) displays a random scatter around the zero line, confirming linearity and homoscedasticity.}
        \label{fig:diagnostics}
\end{figure}

All three diagnostics satisfy the necessary assumptions. The residuals closely align with the reference line in Figure \ref{fig:qq_plot}, confirming normally distributed error terms. This is further substantiated by the symmetric bell-curve in Figure \ref{fig:residuals_histogram}. Finally, the random dispersion around the zero line in Figure \ref{fig:residuals_vs_fitted} confirms the linearity of the relationship, indicating a robust baseline regression model.
\\
\\
\subsection{Limitations}
Unfortunately there are also limitations. There are missing variables such as expectations for unemployment and inflation. Monetary and fiscal policy could also alter the Phillips Curve. There could be omitted variable bias(OVB), OVB is when there are variables that are not in a model but affect the models. Endogeneity is when two variables effect each other and it is impossible to tell what variable affects the other. Endogeneity is present in this paper because of the core relationship between inflation and unemployment that is getting studied..The economy also structurally changes over time which could affect the Phillips Curve. This is also only a simplified model of the economy, this model doesn't take in enough factors to better represent the economy. The trade to GDP doesn't take into account re-exports, re-exports could artificially boosting the trade to GDP number. Using a dummy variable could also hide certain long-term trend affecting the data. Using the dates from 2001-2026 could be skewed due to the 3 recessions and the inflation spike in the aftermath of COVID-19. Also in the second model, the addition of lagged inflation proved that the break was not on the cyclical level, meaning that is it possible that no break occurred.

\section{Discussion}

Looking at the data paints a picture of a Phillips Curve that has revived in certain periods, but is still active despite the theories about globalization. The model also reveals how China joining the WTO and possibly the dot-com crash, lowered baseline inflation over the next 20 years, but didn't affect the Phillips Curve. Though, by adding lagged inflation is can be observed that the important variables become insignificant. Likely meaning that there was a break just not at the cyclical level. Lagged inflation is also very persistent and means that expectations and consumer decisions are very important for fed policy and inflation control. When putting the inflation gap in the equation the p-values become weaker, and the coefficients become flatter. This is backed up by the Wald test, showing that after 2001 the relationships in the US economy fundamentally shifted.

\section{Conclusion}

In the United States, The accession of China into the WTO, Causes a structural break in the Phillips Curve relationship. The Phillips Curve changes throughout different time "regimes" using different ways to measure inflation. The findings were, China possibly caused a drop in inflation and a strengthening in the Phillips Curve after they joined the WTO. The Phillips Curve is still around but has been flattened and strengthened throughout different time periods likely due to the fact that the Phillips Curve was not cyclical. This could be interpreted as firms offshoring employment to low-cost areas raising unemployment. Firms would also be using cheaper goods lowering inflation. In the two models not using China's rise, inflation was heavily dependent on lagged inflation. The broader macro meaning is when firms can offshore hiring for cheaper costs, inflation declines as international supply rises and domestic stayed steady, because the glut of goods was so great  domestic prices were lowered and domestic unemployment in the non-cyclical sector was raised due to off-shored supply. This research could be continued by looking at the Phillips Curve in different time periods and different countries.

\begin{thebibliography}{9}

\bibitem{fed}
Federal Reserve Economic Data.
"CPIAUCSL, UNRATE, MICH, GDP, IMPGS, EXPGS"
\bibitem{IJCB}
International Journal of Central Banking
"Has Globalization Changed the Phillips Curve? Firm-Level Evidence on the Effect of Activity on Prices"
\bibitem{Science}
Science Direct
"Globalization and the slope of the Phillips curve"
\bibitem{Institute}
Becker Friedman Institute
"Globalization, Inflation Dynamics, and
the Slope of the Phillips Curve"
\end{thebibliography}

\end{document}
